\documentclass[../../../include/open-logic-section]{subfiles}

\begin{document}

\olfileid[hi]{sth}{choice}{hartogs}
\olsection{तुलनीयता और हार्टोग्स का लेम्मा}

वह इसका सकारात्मक पक्ष था। अब नकारात्मक पक्ष देखें। चयन के बिना विषय
\emph{अव्यवस्थित} हो जाता है। क्यों, यह देखने के लिए
\cite{Hartogs1915} का यह सुंदर परिणाम देखें:

\begin{lem}[\emph{$\ZF$ में}]\ollabel{HartogsLemma}
किसी भी समुच्चय $A$ के लिए कोई क्रमसूचक $\alpha$ है ऐसा कि
$\cardnless{\alpha}{A}$।
\end{lem}

\begin{proof}
यदि $B\subseteq A$ और $R\subseteq B^2$, तो
\olref[ord-arithmetic][using-addition]{rankcomputation} से
$\tuple{B,R}\subseteq V_{\setrank{A}+4}$। अतः पृथक्करण का उपयोग करके
विचार करें:
\[
	C = \Setabs{\tuple{B, R} \in V_{\setrank{A}+5}}{B\subseteq A 
	\text{ और $\tuple{B,R}$ एक सुव्यवस्था है}}
\]
प्रतिस्थापन और \olref[ordinals][ordtype]{thmOrdinalRepresentation} का
उपयोग करके यह समुच्चय बनाएँ:
\[
	\alpha = \Setabs{\ordtype{B, R}}{\tuple{B, R} \in C}.
\]
\olref[ordinals][basic]{corordtransitiveord} से $\alpha$ एक क्रमसूचक है,
क्योंकि वह क्रमसूचकों का संक्रमणशील समुच्चय है। आखिर यदि
$\gamma\in\beta\in\alpha$, तो किसी $B\subseteq R$ के लिए
$\beta=\ordtype{B,R}$, जिससे
\olref[ordinals][iso]{wellordinitialsegment} के अनुसार किसी $b\in B$ के
लिए $\gamma=\ordtype{B_b,R_b}$, और इस प्रकार $\gamma\in\alpha$।

विरोधाभास द्वारा मान लें कि कोई !!a{injection}
$f\colon\alpha\to A$ है। तब जहाँ:
\begin{align*}
	B &= \ran{f}\\
	R &= \Setabs{\tuple{f(\alpha), f(\beta)} \in A \times A}{\alpha \in \beta}.
\end{align*}
स्पष्टतः $\alpha=\ordtype{B,R}$ और $\tuple{B,R}\in C$। अतः
$\alpha\in\alpha$, जो विरोधाभास है।
\end{proof}

इससे एक गहन परिणाम निष्पन्न होता है:

\begin{thm}[\emph{$\ZF$ में}]
निम्न दावे तुल्य हैं:
\begin{enumerate}
	\item\ollabel{equivwo} सुव्यवस्था अभिगृहीत
	\item\ollabel{equivcompare} किन्हीं समुच्चयों $A$ और $B$ के लिए या तो
	$\cardle{A}{B}$ अथवा $\cardle{B}{A}$
\end{enumerate}
\end{thm}

\begin{proof}
\emph{\olref{equivwo} $\Rightarrow$ \olref{equivcompare}.} $A$ और $B$ को
नियत करें। \olref{equivwo} का उपयोग करने पर सुव्यवस्थाएँ $\tuple{A,R}$ और
$\tuple{B,S}$ हैं।
\olref[ordinals][ordtype]{thmOrdinalRepresentation} का उपयोग करके
$f\colon\alpha\to\tuple{A,R}$ और $g\colon\beta\to\tuple{B,S}$ को
समाकृतिकताएँ मानें। \olref[sth][ordinals][basic]{ordinalsaresubsets} से या
तो $\alpha\subseteq\beta$ अथवा $\beta\subseteq\alpha$। यदि
$\alpha\subseteq\beta$, तो $\comp{f^{-1}}{g}\colon A\to B$ एक
!!a{injection} है, अतः $\cardle{A}{B}$; इसी प्रकार यदि
$\beta\subseteq\alpha$, तो $\cardle{B}{A}$।

\emph{\olref{equivcompare} $\Rightarrow$ \olref{equivwo}.} $A$ को नियत
करें; \olref{HartogsLemma} से कोई क्रमसूचक $\beta$ है जिसके लिए
$\cardnless{\beta}{A}$। \olref{equivcompare} का उपयोग करने पर हमारे पास
$\cardle{A}{\beta}$ है। अतः कोई !!{injection} $f\colon A\to\beta$ है और
क्रम $\Setabs{\tuple{a,b}\in A\times A}{f(a)\in f(b)}$ परिभाषित करके हम इस
अंतःक्षेप का उपयोग $A$ के अवयवों को सुव्यवस्थित करने में कर सकते हैं।
\end{proof}
\noindent
एक तत्काल परिणाम यह है: यदि सुव्यवस्था विफल होती है, तो कुछ समुच्चय आकार
की दृष्टि से \emph{वास्तव में अतुलनीय} होते हैं। अतः यदि सुव्यवस्था विफल
होती है, तो परिमितेतर कार्डिनल अंकगणित अव्यवस्थित होगा। उदाहरण के लिए हमें
यह विचार छोड़ना होगा कि यदि $A$ और $B$ अनंत हों, तो
$\cardeq{\cardeq{A\disjointsum B}{A\times B}}{M}$, जहाँ $M$, $A$ और $B$
में बड़ा है (\olref[card-arithmetic][simp]{cardplustimesmax} देखें)। समस्या
सरल है: यदि हम $A$ और $B$ के आकार की \emph{तुलना} नहीं कर सकते, तो यह पूछना
निरर्थक है कि कौन बड़ा है।

\end{document}
